\documentclass{elsinta}
\usepackage{url}

% --- Memasukkan Data Dokumen Menggunakan Perintah Baru ---
\projectcode{EL4}
\documentname{\ImplementationTitle}
\documenttitle{Teknik Elektro - Sistem Informasi Tugas Akhir}
\capstonetitle{Teknik Elektro - Sistem Informasi Tugas Akhir}
\shortname{  ELSINTA} 
\documentnumber{TA2526.01.099}
\revisionnumber{03}
\publicationdate{22 September 2025}

\revdatefooter{03/10/2025}
% --- Memasukkan Data Halaman Pengesahan (Halaman 2) ---
\currentpage{2} % Halaman yang sedang dicetak

% Data Tim
\ketuatimnama{Mahasiswa 1}
\ketuatimnim{12345678}

\anggotanamaI{Mahasiswa 2}
\anggotanimI{12345678}

\anggotanamaII{Mahasiswa 3}
\anggotanimII{12345678}


% Data Dosen
\pembimbingInama{Dosen 1}
\pembimbingInip{12345678 123456 1 123}

\pembimbingIInama{Dosen 2}
\pembimbingIInip{12345678 123456 1 123}


% Tanggal Pengesahan
\approvaldate{22 September 2025}

\begin{document}

% Panggil perintah untuk membuat halaman sampul
\coverpage
% Pindah ke Halaman Baru
\newpage

% Membuat Halaman Pengesahan (Halaman 2)
% \approvalpage

% Lanjut ke halaman proposal
\newpage

% ======================================
% Bagian Konten Proposal Dimulai di Sini
% ======================================

% Mengatur ulang nomor halaman dan gaya halaman
\clearpage
\pagenumbering{arabic}
\setcounter{page}{1} 
\pagestyle{myfooter} % Ganti dari 'plain' ke 'myfooter'

% Atur margin standar (Opsional, jika Anda ingin margin isi berbeda dari sampul)
% \newgeometry{left=4cm, right=3cm, top=3cm, bottom=3cm} 
% Catatan: Perintah \newgeometry memerlukan paket 'geometry' yang sudah ada di cls.

\tableofcontents
\newpage
\section*{\docsumary}
Sub-bab ini menyajikan secara kritis hasil nyata dari implementasi solusi atau sistem yang telah dikembangkan. Fokus utamanya adalah membuktikan bahwa tujuan proyek telah tercapai dan kebutuhan pengguna/sistem telah terpenuhi melalui analisis kinerja yang komprehensif.

\newpage
\label{sec:impl_result_analysis}



\section{\SubsystemTitle{A}}

    Bagian ini harus secara \textbf{faktual dan terperinci} menyajikan luaran (\textit{output}) dari tahapan implementasi. Sertakan bukti konkret seperti \textit{screenshot} antarmuka, diagram arsitektur \textit{deployment} akhir, dan data mentah hasil pengujian.
    Jelaskan setiap komponen kunci dan \textbf{demonstrasikan fungsionalitasnya} dengan merujuk pada Gambar dan Tabel yang 
    disertakan.
  

\subsection{\ImplementationProceduresTitle}

Bagian ini menguraikan secara rinci langkah-langkah yang dilakukan untuk menerapkan solusi atau sistem yang diusulkan. Ini harus mencakup semua tahapan pelaksanaan, mulai dari persiapan awal hingga pengujian akhir, dan harus cukup detail sehingga pembaca dapat mereplikasi proses tersebut.
\begin{enumerate}
    \item \textbf{Setup dan Konfigurasi:} Jelaskan lingkungan, perangkat keras, perangkat lunak, dan alat yang digunakan. Sertakan rincian mengenai instalasi, konfigurasi parameter, dan penyiapan data.
    \item \textbf{Langkah-langkah Pelaksanaan:} Sajikan urutan logis dari tindakan yang diambil. Ini mungkin termasuk pengkodean, integrasi modul, pengumpulan data, atau modifikasi desain. Gunakan diagram alir atau daftar berurutan untuk meningkatkan kejelasan.
   
    \item \textbf{Jadwal Pelaksanaan (opsional):} Jika relevan, sertakan garis waktu singkat yang menunjukkan kapan setiap tahap implementasi selesai.
\end{enumerate}

\vspace{0.5cm}

\subsection{\ResultsDiscussionTitle}

\label{sec:results_and_discussion}


\textbf{Results and Discussion} (\textit{Hasil dan Pembahasan})
Sub-bab ini adalah inti dari laporan capstone Anda, di mana Anda menyajikan dan menganalisis temuan dari implementasi dan pengujian Anda.
\begin{enumerate}
    \item \textbf{Penyajian Hasil:}
    \begin{enumerate}
        \item Sajikan semua hasil penting secara jelas dan objektif. Gunakan tabel, grafik, tangkapan layar, dan visualisasi data lainnya yang diberi label dengan benar untuk mendukung temuan Anda.
        \item Hasil harus terkait langsung dengan tujuan proyek dan pertanyaan penelitian yang ditetapkan di Bab Pendahuluan.
        \item Sertakan hasil pengujian fungsionalitas, kinerja, atau validasi model.
    \end{enumerate}
    \item \textbf{Pembahasan Hasil (Discussion):}
    \begin{enumerate}
        \item Interpretasikan hasil yang disajikan. Jelaskan apa arti temuan tersebut dalam konteks masalah proyek Anda.
        \item Bandingkan hasil Anda dengan tujuan awal. Apakah proyek berhasil? Sejauh mana?
        \item Bandingkan hasil Anda dengan karya terkait yang ada (jika relevan).
        \item Jelaskan alasan di balik hasil yang diperoleh, termasuk hasil yang tidak terduga atau anomali.
        \item Hubungkan temuan Anda dengan implikasi praktis dan teoretisnya.
    \end{enumerate}
\end{enumerate}

\vspace{0.5cm}

\subsection{\ProblemsSolutionsTitle}
\label{sec:problems_and_solutions}

\textbf{Problems and Solutions} (\textit{Masalah dan Solusi})
Bagian ini menunjukkan kemampuan Anda untuk berpikir kritis dan mengatasi kesulitan yang muncul selama proses proyek. Ini meningkatkan kredibilitas dan transparansi penelitian Anda.
\begin{enumerate}
    \item \textbf{Identifikasi Masalah:} Jelaskan hambatan, tantangan teknis, kendala sumber daya, atau kesulitan metodologis signifikan yang dihadapi selama penelitian, implementasi, atau pengujian.
    \item \textbf{Analisis Masalah:} Berikan analisis singkat tentang sifat dan dampak masalah tersebut terhadap proyek.
    \item \textbf{Solusi yang Diimplementasikan:} Uraikan secara detail langkah-langkah atau perubahan desain yang Anda lakukan untuk mengatasi setiap masalah.
    \item \textbf{Dampak Solusi:} Jelaskan bagaimana solusi yang diterapkan berhasil (atau gagal) dalam mengatasi masalah, dan bagaimana hal tersebut memengaruhi hasil akhir proyek.
    \item \textbf{Pembelajaran (Lessons Learned):} Singgung secara ringkas pelajaran penting yang dipetik dari proses pemecahan masalah.
\end{enumerate}

\section{\SubsystemTitle{B}}
Struktur penulisan sama dengan Subsistem A
\section{\SubsystemTitle{C}}

Struktur penulisan sama dengan Subsistem A
   

\section{\SystemIntegrationTitle}
Struktur penulisan sama dengan Subsistem A



\textbf{System Integration} (\textit{Integrasi Sistem})
Bagian ini menjelaskan proses penggabungan berbagai komponen, modul, atau subsistem yang telah dikembangkan secara terpisah menjadi satu sistem yang kohesif dan berfungsi penuh. Integrasi adalah langkah krusial untuk memastikan bahwa semua bagian bekerja sama sesuai tujuan proyek.
\begin{enumerate}
    \item \textbf{Strategi Integrasi:} Jelaskan pendekatan yang digunakan (misalnya, \textit{Big Bang}, \textit{Incremental}, atau \textit{Sandwich}). Berikan justifikasi mengapa strategi tersebut dipilih.
    \item \textbf{Deskripsi Komponen:} Identifikasi dan jelaskan setiap modul utama atau sistem eksternal yang diintegrasikan.
    \item \textbf{Proses Integrasi:} Uraikan langkah-langkah teknis yang dilakukan untuk menghubungkan komponen, termasuk detail tentang antarmuka (\textit{interface}), API (\textit{Application Programming Interface}), atau protokol komunikasi yang digunakan.
    \item \textbf{Pengujian Integrasi:} Jelaskan bagaimana pengujian dilakukan untuk memverifikasi bahwa modul-modul yang terintegrasi dapat bertukar data dan berfungsi secara harmonis tanpa konflik. Sertakan kasus uji utama yang berfokus pada titik interaksi.
    \item \textbf{Masalah dan Penanganan Integrasi:} Dokumentasikan setiap tantangan atau kegagalan yang terjadi selama proses integrasi dan solusi yang diterapkan untuk mengatasinya.
\end{enumerate}

\vspace{0.5cm}

\section{\UserAcceptanceTitle}
\textbf{User Acceptance} (\textit{Penerimaan Pengguna})
Bagian ini mendokumentasikan proses pengujian akhir yang dilakukan oleh pengguna akhir atau klien proyek untuk memverifikasi bahwa sistem memenuhi kebutuhan bisnis atau fungsionalitas yang disepakati. Ini adalah langkah validasi kritis sebelum proyek dianggap selesai.
\begin{enumerate}
    \item \textbf{Tujuan Uji Penerimaan Pengguna (UAT):} Nyatakan secara jelas tujuan dari pengujian ini, yang biasanya adalah untuk memverifikasi bahwa sistem bekerja dalam skenario dunia nyata dan memenuhi spesifikasi persyaratan awal.
    \item \textbf{Metodologi dan Lingkungan UAT:} Jelaskan siapa yang berpartisipasi (pengguna), kapan pengujian dilakukan, dan di lingkungan mana (misalnya, \textit{staging} atau \textit{production}).
    \item \textbf{Skenario Uji (Test Cases):} Sajikan skenario atau kasus uji utama yang dijalankan oleh pengguna, yang harus didasarkan pada alur kerja (\textit{workflow}) bisnis atau kebutuhan operasional sehari-hari.
    \item \textbf{Hasil UAT:} Dokumentasikan hasil dari pengujian ini, termasuk tingkat keberhasilan (\textit{pass/fail rate}) untuk setiap skenario utama. Gunakan tabel atau ringkasan statistik.
    \item \textbf{Kesimpulan Penerimaan:} Berikan pernyataan formal mengenai penerimaan sistem oleh pengguna. Apakah sistem diterima, diterima dengan syarat, atau ditolak? Jelaskan alasan di balik keputusan tersebut dan tindak lanjut (perbaikan) yang diperlukan jika ada.
\end{enumerate}


\section{\ReferencesTitle}
Daftar pustaka ditulis sesuai standar IEEE, sebaiknya gunakan Citation tools seperti Mendeley, Zotero dll.  
Contoh format IEEE:  
\renewcommand{\refname}{}      % Hapus judul default "Pustaka"
\bibliographystyle{IEEEtran} % Menggunakan style sitasi IEEE
\bibliography{pustaka}    % Memanggil file pustaka.bib

\section{\AbbrevTitle}

% longtable tidak menggunakan lingkungan table[] di sekitarnya.
% Perintah \begin{longtable} langsung menggantikan \begin{table}\begin{tabular}
\begin{longtable}{|p{4cm}|p{11cm}|} % Lebar kolom ditetapkan agar teks wrap
    \caption{Daftar Singkatan (Abbreviations)} \label{tab:abbreviations4} \\
    \hline
    \textbf{Abbreviation} & \textbf{Definition} \\
    \hline
    \endhead % Mengakhiri baris header yang akan diulang di setiap halaman
    
    \hline 
    \multicolumn{2}{|r|}{Lanjutan di halaman berikutnya...} \\ 
    \endfoot % Footer untuk semua halaman kecuali yang terakhir
    
    \hline
    \endlastfoot % Footer untuk halaman terakhir (kosong)

    % --- Isi Tabel ---
    IoT & Internet of Things \\
    \hline
    KPI & Key Performance Indicator \\
    \hline
    ROI & Return on Investment \\
    \hline
    WBS & Work Breakdown Structure \\
    \hline
    SWOT & Strengths, Weaknesses, Opportunities, Threats \\
    \hline
    BMC & Business Model Canvas \\
    \hline
    IEEE & Institute of Electrical and Electronics Engineers \\
    \hline
    ISO & International Organization for Standardization \\
    \hline
    G-MoCo & Gas Quality Monitoring and Controlling (Nama Proyek Anda) \\
    \hline
    CO & Carbon Monoxide \\
    \hline
    CH$_4$ & Methane (atau LPG) \\
    \hline
    PM2.5 & Particulate Matter 2.5 \\
    \hline
    
    % Tambahkan lebih banyak baris di sini untuk memicu pemisahan halaman
    % ... (Tambahkan entri lain jika tabel perlu memanjang)

\end{longtable}
\end{document}